%======================================================================
%  08 参考答案与解析（仅教师版）
%======================================================================
\clearpage
\section*{参考答案与解析}
\addcontentsline{toc}{section}{参考答案与解析}

\newcommand{\ans}[2]{\item \jxlabel{答案}#1\quad\jxlabel{解析}#2}

\subsection*{A组　基础巩固}
\begin{enumerate}[label=\arabic*．,leftmargin=1.8em,itemsep=4pt]
\ans{D}{管道高频焊机利用线圈的交变磁场在圆管中产生涡流，靠涡流的热效应熔化接缝，利用的是电磁感应规律，其发现者是法拉第。库仑发现了点电荷间相互作用的规律，霍尔发现了霍尔效应，洛伦兹提出了磁场对运动电荷的作用力。}

\ans{B}{电磁炉、真空冶炼炉利用涡流的热效应，金属探测器利用涡流的磁效应，都属于涡流的利用；变压器铁芯用相互绝缘的硅钢片叠成是为了减小涡流损耗，属于防止涡流的危害。}

\ans{B}{电磁炉线圈中通高频交变电流，产生的交变磁场在锅底中产生涡流，锅底自身发热来加热食物，A错误、B正确；恒定电流产生恒定磁场，锅底中不会产生涡流，C错误；陶瓷是绝缘体，不能产生涡流，无法用电磁炉加热，D错误。}

\ans{AC}{安检门线圈中通变化的电流，产生变化的磁场，金属物品中产生涡流，涡流的磁场反过来影响线圈中的电流而报警，A正确、B错误；塑料是绝缘体，不能产生涡流，C正确；铝、铜等金属不能被磁化，同样能被安检门检测出来，可见安检门不是利用“金属被磁化”的原理，D错误。}

\ans{AD}{涡流本质上是感应电流，遵循法拉第电磁感应定律和楞次定律，A正确；块状金属静止在恒定磁场中，磁通量不变，不产生涡流，B错误；涡流有磁效应，金属探测器就是利用涡流的磁效应工作的，C错误；频率越高，磁通量变化率越大，感应电动势越大，涡流越大，D正确。}

\ans{A}{由安培定则，螺线管中电流从上往下看沿逆时针方向时，铝板处的磁场竖直向上；电流增大，穿过铝板的磁通量增大，由楞次定律，涡流的磁场竖直向下，所以从上往下看涡流沿顺时针方向，与螺线管中的电流方向相反。方向相反的环形电流相互排斥（“增离减靠”），铝板受到向上的斥力。}

\ans{ABD}{实心铝板在磁场中摆动时产生涡流，受到阻碍其运动的安培力，很快停下，A正确；梳状铝板的缝隙切断了大部分涡流回路，涡流小，阻尼弱，B正确；撤去磁场后没有电磁阻尼，只有很小的空气阻力和摩擦，两摆锤都能摆动较长时间，C错误；电磁阻尼使摆锤的机械能转化为电能，最终转化为内能，D正确。}

\ans{C}{磁体下落时铜管中产生涡流，由“来拒去留”，磁体受到向上的电磁阻力，阻力随速度增大而增大，由 $a=g-\dfrac Fm$ 可知加速度逐渐减小；当阻力增大到等于重力时，磁体匀速下落。下落过程中有机械能转化为内能，机械能不守恒。}

\ans{A}{闭合开关瞬间，线圈中电流从无到有，穿过铝环的磁通量增大，铝环中产生感应电流，由“增离减靠”可知铝环受到向上的斥力而跳起，A正确；电流稳定后磁通量不变，铝环中没有感应电流，不再受力，铝环会落回，B错误；有缺口的铝环不能构成闭合回路，没有感应电流，不受安培力，C错误；电源正、负极对调后，闭合开关瞬间穿过铝环的磁通量仍然增大，铝环仍受斥力而跳起，D错误。}
\end{enumerate}

\subsection*{B组　能力提升}
\begin{enumerate}[label=\arabic*．,leftmargin=1.8em,itemsep=4pt,start=10]
\ans{A}{电子沿逆时针方向（俯视）运动，要使洛伦兹力指向圆心，磁场必须竖直向上，即 $B>0$，对应 $0\sim\dfrac T2$；要使电子加速，电子所受电场力应沿逆时针方向，即感生电场沿顺时针方向，由楞次定律，竖直向上的磁场必须增强，对应 $0\sim\dfrac T4$。在 $\dfrac T4\sim\dfrac T2$ 内磁场减弱，电子被减速；在 $\dfrac T2\sim T$ 内磁场向下，洛伦兹力背离圆心，电子不能沿该轨道运动。}

\ans{ACD}{甲中铜盘转动时，盘上不断有部分进入、离开磁场，这些部分的磁通量变化，产生涡流，涡流受到的安培力阻碍铜盘转动，铜盘很快停下，减少的动能转化为内能，A、D正确。乙中整个铜盘处在匀强磁场中，转动时穿过盘上任一闭合回路的磁通量都不变，不产生涡流，没有电磁阻尼，铜盘（理想情况下）将一直转动，B错误；但铜盘的每条半径都在切割磁感线，产生沿半径方向的感应电动势，圆心与边缘之间存在电势差（这正是法拉第圆盘发电机的原理），C正确。}

\ans{AC}{磁铁靠近铝环正上方时，穿过铝环的磁通量增大，由“来拒”，铝环受到斜向右下方的斥力，其水平分力向右，竖直分力向下，铝环对桌面的压力大于重力，C正确；磁铁离开时，穿过铝环的磁通量减小，由“去留”，铝环受到斜向右上方的引力，水平分力仍向右，竖直分力向上，压力小于重力，D错误。整个过程中铝环所受水平方向的力始终向右，铝环向右运动（电磁驱动），A正确、B错误。}

\item \jxlabel{答案}（1）\qty{0.8}{V}，\qty{4}{m/s^2}　（2）\qty{6.25}{m}，\qty{2.5}{C}　（3）\qty{1.0}{J}　（4）\qty{3.75}{m}

\jxlabel{解析}（1）$E=Blv_0=0.4\times0.5\times\qty{5}{V}=\qty{1.0}{V}$，$I=\dfrac{E}{R+r}=\qty{2}{A}$，$ab$ 两端电压为路端电压 $U=IR=\qty{0.8}{V}$。安培力 $F=BIl=\qty{0.4}{N}$，$a=\dfrac Fm=\qty{4}{m/s^2}$。

（2）由动量定理 $-Blq=0-mv_0$，得 $q=\dfrac{mv_0}{Bl}=\dfrac{0.1\times5}{0.4\times0.5}\,\unit{C}=\qty{2.5}{C}$；又 $q=\dfrac{\dd\Phi}{R+r}=\dfrac{Blx}{R+r}$，得 $x=\dfrac{q(R+r)}{Bl}=\qty{6.25}{m}$。

（3）回路中产生的总焦耳热 $Q=\dfrac12mv_0^2=\qty{1.25}{J}$，$Q_R=\dfrac{R}{R+r}Q=\qty{1.0}{J}$。

（4）由 $v=v_0-\dfrac{B^2l^2}{m(R+r)}x'$ 得 $x'=\dfrac{(v_0-v)m(R+r)}{B^2l^2}=\dfrac{3\times0.1\times0.5}{0.04}\,\unit{m}=\qty{3.75}{m}$。

\item \jxlabel{答案}（1）$\dfrac{kv_0}{2m}$　（2）$\dfrac{mv_0}{2k}$　（3）$\dfrac{mv_0}{2k}+\dfrac{mv_0^2}{8f}$　（4）见解析

\jxlabel{解析}（1）$F=k\cdot\dfrac{v_0}{2}$，$a=\dfrac Fm=\dfrac{kv_0}{2m}$。

（2）由动量定理 $-k\sum v\,\dd t=m\cdot\dfrac{v_0}{2}-mv_0$，而 $\sum v\,\dd t=x_1$，得 $x_1=\dfrac{mv_0}{2k}$。

（3）机械制动阶段加速度大小 $a'=\dfrac fm$，滑行距离 $x_2=\dfrac{(v_0/2)^2}{2a'}=\dfrac{mv_0^2}{8f}$，总距离 $x=x_1+x_2=\dfrac{mv_0}{2k}+\dfrac{mv_0^2}{8f}$。

（4）涡流制动力 $F=kv$ 随速度减小而减小，速度越小，制动力和加速度越小，速度减小得越来越慢，很难在较短时间内停车；而且速度为零时制动力也为零，涡流制动不能使列车保持静止。所以停车前的最后阶段要靠机械制动。
\end{enumerate}

\subsection*{C组　拓展创新}
\begin{enumerate}[label=\arabic*．,leftmargin=1.8em,itemsep=4pt,start=15]
\item \jxlabel{答案}（1）$\dfrac{k\pi r^2}{R_0}$，逆时针方向　（2）$\dfrac{k\pi R^2}{R_0}$，逆时针方向　（3）$\dfrac{kr}{2}$，$\dfrac{kR}{4}$

\jxlabel{解析}（1）$E_1=\dfrac{\dd B}{\dd t}\cdot\pi r^2=k\pi r^2$，$I_1=\dfrac{k\pi r^2}{R_0}$。磁场垂直纸面向里且增强，由楞次定律，感应电流的磁场垂直纸面向外，由安培定则，感应电流沿逆时针方向。

（2）半径为 $2R$ 的圆环中，有磁场的面积仍为 $\pi R^2$，$E_2=k\pi R^2$，$I_2=\dfrac{k\pi R^2}{R_0}$，方向仍为逆时针。

（3）由对称性，感生电场线是同心圆，同一圆周上场强大小相等。感生电场力移动单位正电荷绕行一周做的功等于感生电动势：

$r<R$ 处：$E_{\text{感}}\cdot2\pi r=k\pi r^2$，得 $E_{\text{感}}=\dfrac{kr}{2}$；

$2R$ 处：$E_{\text{感}}'\cdot2\pi(2R)=k\pi R^2$，得 $E_{\text{感}}'=\dfrac{kR}{4}$。可见磁场区域外也有感生电场。

\item \jxlabel{答案}（1）$eK$　（2）$\dfrac{K}{2\pi r}$　（3）$NeK$　（4）$\dfrac{K}{2\pi r^2}$

\jxlabel{解析}（1）轨道所围回路中的感生电动势 $E=\dfrac{\dd\Phi}{\dd t}=K$，电子绕行一周，感生电场力做功 $W=eK$。

（2）$E_{\text{感}}\cdot2\pi r=K$，得 $E_{\text{感}}=\dfrac{K}{2\pi r}$。

（3）洛伦兹力不做功，由动能定理 $\dd E_{\mathrm k}=NW=NeK$。

（4）径向：$evB_r=m\dfrac{v^2}{r}$，得 $mv=eB_rr$，所以 $\dd(mv)=er\,\dd B_r$。切向：由动量定理 $\dd(mv)=eE_{\text{感}}\dd t=\dfrac{eK}{2\pi r}\dd t$。联立得
\[\frac{\dd B_r}{\dd t}=\frac{K}{2\pi r^2}.\]
（轨道所围面积内平均磁感应强度的变化率为 $\dfrac{K}{\pi r^2}$，可见 $\dfrac{\dd B_r}{\dd t}$ 恰为它的一半，这就是“2:1”条件。）

\item \jxlabel{答案}（1）49\%　（2）64 倍，$\dfrac18$　（3）不合理，见解析

\jxlabel{解析}（1）其他条件不变时，单位体积涡流损耗 $\dfrac PV\propto d^2$，铁芯总体积不变，所以总损耗变为原来的 $\left(\dfrac{0.35}{0.5}\right)^2=0.49$，即 49\%。

（2）$\dfrac PV\propto f^2$，频率变为原来的 8 倍，损耗变为原来的 $8^2=64$ 倍；要使损耗不变，需 $fd$ 不变，所以硅钢片厚度应变为原来的 $\dfrac18$。

（3）不合理。变压器铁芯的主要作用是增强磁场、约束磁感线、减少漏磁，必须用导磁性能好的铁磁性材料；普通陶瓷不是铁磁性材料，不能增强磁场，用它做“铁芯”与没有铁芯差不多，变压器的效率会大大降低。正确的思路是：选用\textbf{既能导磁、电阻率又大}的材料，如硅钢片（并叠成绝缘薄片），在高频时则用铁氧体——一种电阻率很大的“磁性陶瓷”。
\end{enumerate}
